\clearpage\begingroup\NoSourceHead\SourceLeaf{507}
\renewcommand{\thefootnote}{\arabic{footnote})}\setcounter{footnote}{72}
\ArithmeticHead{50}{L}{BOLZA, GAUSS UND DIE VARIATIONSRECHNUNG.}
\noindent\hspace*{\parindent}27. Während es sich bei den Untersuchungen, über die wir in den vorangehenden
Nummern berichtet haben, darum handelte, die \textsc{Gauss}schen Beweise
zu vereinfachen oder zu verallgemeinern, haben wir nun zum Schluss noch
einige \textls[100]{Anwendungen} zu besprechen, die von der \textsc{Gauss}schen Formel für
die Variation des Flächeninhalts \textls[100]{auf Scharen von Minimalflächen} gemacht
worden sind.

Da ist zunächst der folgende Satz zu erwähnen, den \textsc{Lamarle}\footnote{\textsc{Lamarle}, loc. cit. 67), S. 576.} (1863)
aus der Formel (64) abgeleitet hat: »Es sei eine einparametrige Schar $A$
von Minimalflächen gegeben; auf einer derselben ziehe man eine geschlossene
Kurve und lege durch sie eine Fläche $B$, die sämtliche Flächen der Schar
orthogonal schneidet. Alsdann haben die Flächenstücke, die durch die röhrenförmige
Fläche $B$ aus den verschiedenen Flächen der Schar $A$ ausgeschnitten
werden, alle denselben Flächeninhalt«.

Noch wichtiger für die Variationsrechnung ist ein hiermit nahe verwandter
Satz von H. A. \textsc{Schwarz}\footnote{\textsc{Schwarz}, loc. cit. 72), S. 224—227. Wir haben uns in der Bezeichnung möglichst nahe an die
Bezeichnung von \textsc{Gauss} angeschlossen.} (1885). \textsc{Schwarz} geht von der Bemerkung
aus, dass der Ausdruck (45) für $\delta U$ in dem besonderen Fall, wo $U$ eine Minimalfläche
ist, sich auf das Linienintegral reduziert und sich daher unter Benutzung
von (46) schreiben lässt
\begin{equation*}\tag{71}
\delta U=-\int\delta p\mathpunct{.}\cos i\mathpunct{.}dP,
\end{equation*}
wenn man mit $\delta p$ die Projektion des Variationsvektors auf die Richtung 8 bezeichnet.
\textsc{Schwarz} betrachtet nun eine einparametrige Schar von Minimalflächenstücken
$U(\varepsilon)$, deren Begrenzungskurven $P(\varepsilon)$ sämtlich auf der Gefässwand
liegen sollen. Ferner wird angenommen, dass keine zwei Minimalflächenstücke
einen gemeinsamen Punkt haben sollen, wenn $\varepsilon$ auf ein bestimmtes
Intervall, etwa $(0,1)$ beschränkt wird. Den Werten $\varepsilon=0$ und $\varepsilon=1$ mögen
die Flächen $U_0$ und $U$ entsprechen. Während $\varepsilon$ von 0 bis 1 wächst, überdeckt
die Schar der Begrenzungskurven $P(\varepsilon)$ einen gürtelförmigen Streifen $F$
der Gefässwand $T$. Es wird jetzt der aufsteigenden Reihe
\[
\varepsilon=0,\ \varepsilon',\ \varepsilon'',\ .\, .\, .,\ \varepsilon^{(n)},\ 1
\]
entsprechend eine Reihe von Minimalflächenstücken
\clearpage\SourceLeaf{508}
\ArithmeticHead{51}{R}{WIRKUNGEN DER GAUSSSCHEN ARBEIT.}
\begin{equation*}\tag{72}
U_0,\ U',\ U'',\ .\, .\, .\ \ U^{(n)},\ U
\end{equation*}
aus der Schar herausgegriffen. Für die Differenz der Flächeninhalte zweier
aufeinanderfolgender Flächenstücke der Reihe gilt dann angenähert die Formel
(71). Gleichzeitig wird durch die Begrenzungskurven der Flächenstücke
(72) zusammen mit einer Anzahl von orthogonalen Trajektorien der Schar
$P(\varepsilon)$ der Flächenstreifen $F$ in Elemente geteilt; der Flächeninhalt $dF$ eines
solchen Elementes ist gegeben durch
\[
dF=\delta p\mathpunct{.}dP.
\]

Daraus folgt durch Summation und Grenzübergang
\begin{equation*}\tag{73}
U-U_0=-\int\cos i(\varepsilon)\,dF,
\end{equation*}
wobei das Oberflächenintegral über den Streifen $F$ zu erstrecken ist, und $i(\varepsilon)$
den »Randwinkel« für die durch den Ort des Elementes $dF$ gehenden Minimalfläche
$U(\varepsilon)$ bedeutet.

Wenn insbesondere die Fläche $U_0$ auf einen Punkt zusammenschrumpft
und dementsprechend der Streifen $F$ in die Fläche $T$ übergeht, so geht die
Gleichung (73) über in
\[
U=-\int\cos i(\varepsilon)\,dT,
\]
erstreckt über die Fläche $T$; daraus folgt aber, wenn man noch: $\omega=\pi-i(\varepsilon)$
setzt, der \textsc{Schwarz}sche Fundamentalsatz
\[
T-U=\int(1-\cos\omega)\,dT.
\]

Dieser Satz entspricht für das Problem der Minimalflächen dem \textsc{Weierstrass}schen
Fundamentalsatz über die Darstellung der vollständigen Variation
eines bestimmten Integrals durch die $E$-Funktion. \textsc{Schwarz} hat auch einen
rein analytischen Beweis seines Satzes gegeben, der auf der in der vorigen
Nummer erwähnten Formel (69) beruht.
\par\begin{center}\rule{24mm}{0.3pt}\end{center}
\clearpage\endgroup
